\subsection{Влияние динамического перекрытия поля зрения на качество одометрии}
\label{sec:1-3-dynamic-overlap}

Для лидарной одометрии особенно сложны дорожные сцены, в которых значительная часть поля зрения временно перекрывается крупным динамическим объектом. Практически такая ситуация возникает, например, при движении рядом с грузовиком или при его обгоне по соседней полосе. В этом случае часть лучей лидара начинает попадать в кузов соседнего автомобиля, а часть стационарной сцены за ним перестаёт наблюдаться. Тем самым входные данные изменяются не только из-за собственного движения беспилотного автомобиля, но и из-за того, что наблюдаемая геометрия сцены становится неполной.

Для алгоритмов регистрации облаков точек это принципиально важно, поскольку их работа опирается на наличие достаточного числа корректных соответствий между соседними кадрами. При динамическом перекрытии поля зрения одновременно проявляются несколько механизмов деградации. Во-первых, уменьшается доля наблюдаемых стационарных ориентиров, по которым можно надёжно оценивать относительное движение. Во-вторых, возрастает доля точек, которые принадлежат самому динамическому объекту и потому не подчиняются единой модели собственного движения автомобиля. В-третьих, ухудшается геометрическая обусловленность задачи регистрации: даже если часть статической сцены остаётся видимой, она может оказаться менее информативной, например состоять преимущественно из удалённых, слабо структурированных или почти плоских участков.

Важно, что проблема не сводится только к появлению выбросов. Даже если часть явно неудачных соответствий удаётся отбросить пороговой фильтрацией или ослабить их влияние с помощью весов, дающих устойчивость к выбросам, это ещё не гарантирует устойчивой оценки движения. При сильном перекрытии поля зрения может оказаться, что после удаления сомнительных пар остаётся слишком мало информативных соответствий, чтобы надёжно восстановить относительное преобразование. Иначе говоря, алгоритм сталкивается не только с загрязнением данных, но и с дефицитом полезной геометрической информации.

Для классического ICP это означает рост риска локальной деградации оценки движения. Даже если алгоритм в среднем работает устойчиво на обычных последовательностях, при существенном перекрытии поля зрения число полезных соответствий может оказаться недостаточным для точной регистрации, а ближайшие соседи начинают хуже отражать реальную структуру статической сцены. В результате оценка преобразования сильнее зависит от начального приближения, конкретной конфигурации наблюдаемых объектов и доли динамических точек в кадре.

Таким образом, анализ существующих подходов показывает, что классические методы регистрации облаков точек, и прежде всего ICP в базовой постановке, предполагают наличие достаточно большого перекрытия статической части сцены и ограниченного числа ложных соответствий. В условиях динамического перекрытия поля зрения эти предпосылки нарушаются, что приводит к деградации оценки относительного движения. Поэтому в рассматриваемой задаче недостаточно ограничиться только базовой схемой регистрации соседних облаков точек. Требуется воспроизводимая экспериментальная постановка, в которой перекрытие поля зрения можно контролируемо моделировать, а также устойчивый алгоритм, уменьшающий влияние неудачных соответствий и использующий априорную модель движения.
